% !PRINTED 195
\weberpage{195}
Darin liegt auch der Beweis, dass, wenn $D$ von Null verschieden ist, die Normalform durch lineare Transformation hergestellt werden kann. Denn unter dieser Voraussetzung zerfällt $H$ in zwei von einander verschiedene lineare Factoren, und wenn man diese, von constanten Factoren abgesehen, für die neuen Variablen $\xi$, $\eta$ einer linearen Transformation wählt, so wird $A_0' = 0$, $A_2' = 0$; es folgt aber aus den identischen Relationen
\[
3\,a_3\,A_0 + a_1\,A_2 = a_2\,A_1,\quad a_2\,A_0 + 3\,a_0\,A_2 = a_1\,A_1,
\]
wenn nicht zugleich $A_1' = 0$, was durch das nicht verschwindende $D$ ausgeschlossen ist, $a_1' = 0$, $a_2' = 0$, d.\,h. die transformirte Form von $f$ enthält nur die Cuben von $\xi$ und $\eta$.

Erhebt man die erste Gleichung (10) in den Cubus, und eliminirt $\xi^3$, $\eta^3$, $r^6$ mittelst (11) und der zweiten Gleichung (10), so erhält man zwischen den Covarianten folgende identische Relation
\weq{12}{4\,H^3 + Q^2 + 27\,D\,f^2 = 0.}

Um die Transformation in die Normalform auszuführen, d.\,h. die Functionen $\xi$, $\eta$ wirklich zu finden, zerlegt man die Function $H$ in ihre linearen Factoren
% [sic] Im Druck steht im zweiten Factor A_2 y. Rechnerisch muss dort wie im ersten Factor A_1 y stehen, denn A_1^2-4A_0A_2=-3D. Der Druck wird unverändert bewahrt.
\[
4\,A_0\,H = \bigl(2\,A_0\,x + A_1\,y + \sqrt{-3\,D}\,y\bigr)\,\bigl(2\,A_0\,x + A_2\,y - \sqrt{-3\,D}\,y\bigr).
\]
Dann unterscheiden sich $\xi$, $\eta$ von diesen beiden Factoren von $H$ nur um je einen constanten Factor. Wir setzen also, wenn wir zwei constante Factoren mit $h$, $k$ bezeichnen,
\weq{13}{
\begin{aligned}
2\,\xi &= h\,\bigl(2\,A_0\,x + A_1\,y - \sqrt{-3\,D}\,y\bigr) \\
2\,\eta &= k\,\bigl(2\,A_0\,x + A_1\,y + \sqrt{-3\,D}\,y\bigr),
\end{aligned}
}
woraus sich durch Multiplication mit Rücksicht auf die beiden ersten Gleichungen (10) ergiebt
\weq{14}{3\,h\,k\,A_0\,\sqrt[3]{-D} = 1,}
und die Gleichungen (11) geben durch Vergleichung der Co\"efficienten von $x^3$
\weq{15}{
\begin{aligned}
6\,h^3\,A_0^3\,\sqrt{-3\,D} &= +\,q_0 + 3\,\sqrt{-3\,D}\,a_0 \\
6\,k^3\,A_0^3\,\sqrt{-3\,D} &= -\,q_0 + 3\,\sqrt{-3\,D}\,a_0,
\end{aligned}
}
wo $q_0$ der Co\"efficient von $x^3$ in $Q$, also
\weq{16}{q_0 = 27\,a_0^2\,a_3 - 9\,a_0\,a_1\,a_2 + 2\,a_1^3}
ist.
